\subsection{平坦性的局部化}

命题 13. 设 S 是环 A 的一个乘性子集，M 是一个 A-模。若 M 是平坦的（相应地，忠实平坦的），则 \(S^{-1}M\) 是平坦（相应地，忠实平坦）\(S^{-1}A\)-模，并且是平坦 A-模。

由于 \(S^{-1}M = M \otimes_{A} S^{-1}A\)，第一个断言由第 I 章 §2 第 7 节命题 8 的推论 2（相应地，第 I 章 §3 第 3 节命题 5）得出；此外，\(S^{-1}A\) 是平坦 A-模（§2 第 4 节定理 1）；因此若 M 是平坦 A-模，则凭第 I 章 §2 第 7 节命题 8 的推论 3，\(S^{-1}M\) 也是。

注. 若 N 是一个 \(S^{-1}A\)-模，则 \(S^{-1}N\) 等同于 N，因而这等价于说 N 是平坦 \(S^{-1}A\)-模或平坦 A-模。

命题 14. 设 A 是一个环，B 是一个交换 A-代数，T 是 B 的一个乘性子集。若 N 是一个作为 A-模时平坦的 B-模，则 \(T^{-1}N\) 是平坦 A-模。

\(T^{-1}N = T^{-1}B \otimes_{B} N\)；于是该命题由第 I 章 §2 第 7 节命题 8 得出，其中 A 换成 B，B 换成 A，E 换成 \(T^{-1}B\)，F 换成 N。

命题 15. 设 A、B 是两个环，\(\phi: A \to B\) 是一个同态，N 是一个 B-模。下列性质等价：

\begin{enumerate}
\def\labelenumi{(\alph{enumi})}
\item
  N 是平坦 A-模。
\item
  对每个 \(\mathfrak{n}\)（即 \(\mathbf{B}\) 的极大理想），\(\mathbf{N}_{\mathfrak{n}}\) 是平坦 A-模。
\item
  对每个 \(\mathfrak{n}\)（即 \(\mathbf{B}\) 的极大理想），若记 \(\mathfrak{m} = \bar{\phi}^{-1}(\mathfrak{n})\)，则 \(\mathbf{N}_{\mathfrak{n}}\) 是平坦 \(\mathbf{A}_{\mathfrak{m}}\)-模。
\end{enumerate}

对所有 \(a \notin m\)，由 a 诱导的 \(N_{n}\) 的位似是双射，故 \(N_{n}\) 典范地等同于 \((\mathrm{N}_{\mathrm{n}})_{\mathrm{m}}\)，于是 (b) 与 (c) 的等价性由

命题 13 后面的注得出；(a) 蕴含 (b) 这一事实是命题 14 的特殊情形。剩下的要证明 (b) 蕴含 (a)，即若 (b) 成立，则对每个单的 A-模同态 \(u: \mathbf{M} \to \mathbf{M}'\)，同态 \(v = 1 \otimes u: \mathbf{N} \otimes_{\mathbf{A}} \mathbf{M} \to \mathbf{N} \otimes_{\mathbf{A}} \mathbf{M}'\) 是单射。现在，\(v\) 也是 B-模同态，而为使它是单射，必须且只需 \(v_n: (\mathbf{N} \otimes_{\mathbf{A}} \mathbf{M})_n \to (\mathbf{N} \otimes_{\mathbf{A}} \mathbf{M}')_n\) 对每个 \(n\)（即 \(B\) 的极大理想）都是单射（第 3 节定理 1）。由于

\[
(\mathbf {N} \otimes_ {\mathbf {A}} \mathbf {M}) _ {n} = \mathbf {B} _ {n} \otimes_ {\mathbf {B}} (\mathbf {N} \otimes_ {\mathbf {A}} \mathbf {M}) = \mathbf {N} _ {n} \otimes_ {\mathbf {A}} \mathbf {M},
\]

\(v_{n}\) 正是同态 \(1 \otimes u: N_{n} \otimes_{A} M \to N_{n} \otimes_{A} M'\)，由于按假设 \(N_{n}\) 是平坦 A-模，它是单射的。

推论. 为使一个 A-模 M 是平坦的（相应地，忠实平坦的），必须且只需对 A 的每个极大理想 m，\(M_{m}\) 是平坦（相应地，忠实平坦）\(A_{m}\)-模。

条件的必要性由命题 13 得出。反之，若对 A 的每个极大理想 m，\(M_{m}\) 是平坦 \(A_{m}\)-模，则凭命题 15 应用于 \(\phi\) 是恒同的情形，M 是平坦 A-模。最后，若对一切 m，\(M_{m}\) 是忠实平坦 \(A_{m}\)-模，则 \(mM_{m} = mA_{m}M_{m} \neq M_{m}\)，从而对一切 m 有 \(mM \neq M\)（第 3 节命题 9），这就证明 M 是忠实平坦 A-模（第 I 章 §3 第 1 节命题 1 (d)）。

